\chapter{Annexes}
\label{chap:annexes}

% =====================================================================
%  RAPPEL INSTRUCTION : un diagramme de Gantt FINAL doit figurer en annexe
%  (vue synthétique de la chronologie). Des extraits de code « choisis »
%  sont possibles (jamais des dizaines de pages).
% =====================================================================

\section{Diagramme de Gantt}
Le diagramme ci-dessous restitue la chronologie des missions sur les deux
années d'alternance, le corps du mémoire les présentant quant à lui
regroupées par nature.

% TODO (auteur) : AJUSTER les dates réelles de chaque mission (bornes posées
% ici d'après la trame du mémoire : Free Mobile puis Beemo début 2026,
% framework « étalé sur près d'un an et demi » avec son outillage).
\begin{figure}[H]
\centering
\resizebox{\textwidth}{!}{%
\begin{ganttchart}[
    hgrid, vgrid,
    time slot format=isodate-yearmonth,
    time slot unit=month,
    x unit=0.6cm, y unit chart=0.62cm,
    title label font=\footnotesize,
    bar label font=\small,
    group label font=\small\bfseries,
    bar/.append style={fill=coral!70, draw=coraldeep},
    group/.append style={fill=muted, draw=muted},
    milestone/.append style={fill=coraldeep, draw=coraldeep}
  ]{2024-09}{2026-08}
  \gantttitlecalendar{year} \\
  \gantttitlecalendar{month=shortname} \\
  \ganttgroup{Free Mobile}{2024-09}{2026-01} \\
  \ganttbar{Duplication Grafana}{2024-09}{2024-10} \\
  \ganttbar{Coût de compilation C++}{2024-11}{2024-12} \\
  \ganttbar{Framework de test}{2025-01}{2025-08} \\
  \ganttbar{Extension VS Code (LSP/DAP)}{2025-08}{2025-11} \\
  \ganttbar{Portage forks legacy}{2025-11}{2026-01} \\
  \ganttgroup{Beemo Technologie}{2026-02}{2026-08} \\
  \ganttbar{Rétro-ingénierie DSDK}{2026-02}{2026-04} \\
  \ganttbar{Immuabilité (en cours)}{2026-04}{2026-08}
\end{ganttchart}%
}
\caption{Chronologie des missions de l'alternance.}
\label{fig:gantt}
\end{figure}

\section{Extrait de code choisi : un scénario de test complet}
L'extrait suivant complète le chapitre~\ref{chap:t3} en montrant un scénario
de test du framework dans son intégralité : préparation de l'abonné,
simulation d'un flux d'appel mêlant voix et SMS, puis vérifications par
\emph{snapshot}.

\begin{lstlisting}[caption={Scénario complet de non-régression sur un flux voix + SMS (illustratif).},label={lst:annexe-scenario}]
# --- Preparation -----------------------------------------------------
- declarer_abonne:          # abonne de test, solde initial connu
    msisdn: "33600000001"
    solde: 15.0
    offre: "forfait_2euros"

# --- Flux d'appel simule --------------------------------------------
- appel_voix:               # trames TCP identiques a celles du reseau
    appelant: "33600000001"
    appele:   "33600000002"
    duree: 180              # secondes

- envoi_sms:
    emetteur: "33600000001"
    destinataire: "33600000002"

# --- Verifications (comparees au snapshot de reference) -------------
- verifier_valorisation:    # tickets produits et solde decompte
    msisdn: "33600000001"

- verifier_notifications:   # seuils de consommation eventuels
    msisdn: "33600000001"
\end{lstlisting}
